\subsection{Lorentz geometry and canonical coordinates}
\label{app:rigid-geometry}
Curvature adaptation moves both the point and its tangent space. The purpose of the following identities is to place these moving objects in a common coordinate system while retaining their physical metric. Keeping track of the metric scaling is essential: preserving normalized coordinates does not preserve physical distances or the numerical value of the gradient.
\begin{definition}[State space]
\begin{align}
 J_d&=\operatorname{diag}(-1,I_d),&
 \langle a,b\rangle_L&=a^\top J_db,\\
 H_K^d&=\{x\in\mathbb R^{d+1}:x^\top J_dx=-K^{-1},\ x_0>0\},&
 o_K&=(K^{-1/2},0,\ldots,0),\\
 T_xH_K^d&=\{v:x^\top J_dv=0\},&
 \|v\|_x^2&=v^\top J_dv,\\
 \mathcal Q&=\prod_{g=1}^G(H_1^{d_g})^{n_g},&
 U&=\prod_{g=1}^G[u_{g,-},u_{g,+}],\\
 K_g&=e^{u_g}\in[K_-,K_+],&0&<K_-\le K_+<\infty.
\end{align}
\end{definition}

\begin{lemma}[Positive tangent metric and projection]
For $x=(x_0,x_s)\in H_K^d$,
\begin{align}
 x_0^2&=K^{-1}+\|x_s\|_2^2,&
 v_0&=x_s^\top v_s/x_0\quad(v\in T_xH_K^d),\\
 \|v\|_x^2
 &=\|v_s\|_2^2-(x_s^\top v_s)^2/x_0^2
 \ge\frac{\|v_s\|_2^2}{Kx_0^2}\ge0,\\
 \Pi_x^K a&=a+K\langle x,a\rangle_Lx,&
 (\Pi_x^K)^2&=\Pi_x^K,\\
 \langle x,\Pi_x^Ka\rangle_L&=0,&
 \operatorname{grad}f(x)&=\Pi_x^KJ_d\nabla\widetilde f(x).
\end{align}
\end{lemma}
\begin{proof}
\begin{align}
 (x_s^\top v_s)^2
 &\le\|x_s\|_2^2\|v_s\|_2^2,\\
 1-\|x_s\|_2^2/x_0^2&=(Kx_0^2)^{-1},\\
 \langle x,\Pi_x^Ka\rangle_L
 &=\langle x,a\rangle_L+K\langle x,a\rangle_L(-K^{-1})=0,\\
 \langle\Pi_x^KJ_d\nabla\widetilde f,v\rangle_L
 &=\nabla\widetilde f^\top v
   +K\langle x,J_d\nabla\widetilde f\rangle_L\langle x,v\rangle_L
 =Df(x)[v].
\end{align}
\end{proof}

\begin{proposition}[Connection, curvature and geodesics]
For tangent fields $V,W$ and $v\in T_xH_K^d$,
\begin{align}
 \nabla_VW&=\Pi_x^KD_VW=D_VW-K\langle V,W\rangle_Lx,\\
 R(U,V)W&=-K(\langle V,W\rangle_LU-\langle U,W\rangle_LV),\\
 \operatorname{Sec}(U,V)&=-K,\\
 \gamma_v(s)&=\cosh(\sqrt K\|v\|s)x+
 \frac{\sinh(\sqrt K\|v\|s)}{\sqrt K\|v\|}v,\\
 \gamma_v''(s)&=K\|v\|^2\gamma_v(s),&
 \nabla_{\dot\gamma}\dot\gamma&=0,\\
 \Exp_x^K(v)&=\gamma_v(1),&
 d_K(x,y)&=K^{-1/2}\operatorname{arcosh}(-K\langle x,y\rangle_L).
\end{align}
\end{proposition}
\begin{proof}
\begin{align}
 D_V\langle x,W\rangle_L&=0
 \Longrightarrow\langle x,D_VW\rangle_L=-\langle V,W\rangle_L,\\
 D_VW&=\nabla_VW+K\langle V,W\rangle_Lx,\\
 \Pi_x^K(D_UD_VW-D_VD_UW-D_{[U,V]}W)&=0,\\
 R(U,V)W
 &=-K\langle V,W\rangle_LU+K\langle U,W\rangle_LV,\\
 \langle\gamma_v(s),\gamma_v(s)\rangle_L
 &=-K^{-1}\cosh^2(\sqrt K\|v\|s)
   +K^{-1}\sinh^2(\sqrt K\|v\|s)\\
 &=-K^{-1},\\
 \langle\dot\gamma_v,\dot\gamma_v\rangle_L&=\|v\|^2,\\
 \Pi_{\gamma_v}^K\gamma_v''&=K\|v\|^2\Pi_{\gamma_v}^K\gamma_v=0,\\
 -K\langle x,\gamma_v(s)\rangle_L&=\cosh(\sqrt K\|v\|s).
\end{align}
\end{proof}

\begin{proposition}[Closed-form parallel transport]
\label{prop:rigid-pt}
\begin{align}
 P_{x\to y}^Kv
 &=v+\frac{K\langle y,v\rangle_L}{1-K\langle x,y\rangle_L}(x+y),\\
 \langle y,P_{x\to y}^Kv\rangle_L&=0,&
 \|P_{x\to y}^Kv\|_y&=\|v\|_x,\\
 P_{y\to x}^KP_{x\to y}^K&=I_{T_xH_K}.
\end{align}
\end{proposition}
\begin{proof}
Set $a=\langle y,v\rangle_L$, $b=-K\langle x,y\rangle_L\ge1$,
$c=Ka/(1+b)$. Then
\begin{align}
 \langle y,v+c(x+y)\rangle_L
 &=a-c(1+b)/K=0,\\
 \langle x+y,x+y\rangle_L&=-2(1+b)/K,\\
 \|v+c(x+y)\|_L^2-\|v\|_L^2
 &=2ca-2c^2(1+b)/K=0.
\end{align}
For $y=\gamma_w(s)$, $\lambda=\sqrt K\|w\|$ and
$v=v_\perp+\langle w,v\rangle_Lw/\|w\|^2$,
\begin{align}
 c(s)&=\frac{K\langle w,v\rangle_L}{\lambda}\tanh(\lambda s/2),\\
 P_{x\to\gamma_w(s)}^Kv
 &=v_\perp+\frac{\langle w,v\rangle_L}{\|w\|^2}\dot\gamma_w(s),\\
 \Pi_{\gamma_w(s)}^K\frac d{ds}P_{x\to\gamma_w(s)}^Kv
 &=\frac{\langle w,v\rangle_L}{\|w\|^2}
   \Pi_{\gamma_w(s)}^K\lambda^2\gamma_w(s)=0.
\end{align}
\end{proof}

\begin{proposition}[Homothety and compatible state transport]
\label{prop:homothety}
For $a=\sqrt{K/K'}$, $\Psi_{K\to K'}x=ax$,
$C_{K\to K'}v=av$,
\begin{align}
 \langle ax,ax\rangle_L&=-K'^{-1},&
 \langle ax,av\rangle_L&=0,\\
 \Psi_{K\to K'}^*g_{K'}&=a^2g_K,&
 \|av\|_{ax}&=a\|v\|_x,\\
 d_{K'}(ax,ay)&=a\,d_K(x,y),\\
 \Psi_{K\to K'}\Exp_x^K(v)&=\Exp_{ax}^{K'}(av),\\
 P_{ax\to ay}^{K'}(av)&=aP_{x\to y}^Kv.
\end{align}
\end{proposition}
\begin{proof}
\begin{align}
 a^2/K&=1/K',&\sqrt{K'}\|av\|&=\sqrt K\|v\|,\\
 -K'\langle ax,ay\rangle_L&=-K\langle x,y\rangle_L,\\
 \Exp_{ax}^{K'}(av)
 &=a\cosh(\sqrt K\|v\|)x+
 a\frac{\sinh(\sqrt K\|v\|)}{\sqrt K\|v\|}v,\\
 P_{ax\to ay}^{K'}(av)
 &=av+\frac{K'a^2\langle y,v\rangle_L}{1-K'a^2\langle x,y\rangle_L}a(x+y)
 =aP_{x\to y}^Kv.
\end{align}
\end{proof}

The homothety commutes with geodesic motion and parallel transport, so it provides a consistent identification of the chosen state across curvatures. This compatibility is the property used below. In particular, an old momentum must undergo the same differential as the point; retaining its ambient entries produces the discrepancy quantified after the canonical identities.

\begin{definition}[Canonical trivialization]
\begin{align}\label{eq:canonical}
 q_g&=e^{u_g/2}x_g,&p_g&=e^{u_g/2}m_g,\\
 F(q,u)&=\widetilde{\mathcal J}((e^{-u_g/2}q_g)_g,u),\\
 A(u)&=\operatorname{diag}(e^{u_g}I_{n_gd_g})_{g=1}^G.
\end{align}
\end{definition}
\begin{corollary}[Metric and gradient identities]
\label{cor:canonical-gradient}
\begin{align}\label{eq:equivalent-metrics}
 K_+^{-1}\|v\|_1^2&\le\sum_gK_g^{-1}\|v_g\|_1^2
 \le K_-^{-1}\|v\|_1^2,\\
 \operatorname{grad}_{q_g}F&=K_g^{-1/2}\operatorname{grad}_{x_g}\mathcal J,\\
 \sqrt{K_g}\operatorname{grad}_{x_g}\mathcal J
 &=K_g\operatorname{grad}_{q_g}F,\\
 q_g'&=\sqrt{K_g'}\sqrt{K_g/K_g'}x_g=q_g,\\
 p_g'&=\sqrt{K_g'}\sqrt{K_g/K_g'}m_g=p_g.
\end{align}
\end{corollary}
\begin{proof}
\begin{align}
 D_{q_g}F[v_g]
 &=D_{x_g}\mathcal J[K_g^{-1/2}v_g]\\
 &=\langle K_g^{-1/2}\operatorname{grad}_{x_g}\mathcal J,v_g\rangle_L.
\end{align}
\end{proof}

\begin{proposition}[Omitted transport]
\label{prop:omitted-rigid}
For $u'=u+\delta u$ and $m'=m$,
\begin{align}\label{eq:missing-transport}
 p'&=e^{\delta u/2}p,&
 \|p'-p\|&=|e^{\delta u/2}-1|\|p\|,\\
 \|p'-p\|&\le\tfrac12e^{|\delta u|/2}|\delta u|\|p\|,\\
 p'-p&=\tfrac12\delta u\,p+\tfrac18(\delta u)^2p+O(|\delta u|^3\|p\|),\\
 \prod_{j=s}^{t}e^{(u_{j+1}-u_j)/2}
 &=e^{(u_{t+1}-u_s)/2}\in[\sqrt{K_-/K_+},\sqrt{K_+/K_-}].
\end{align}
\end{proposition}
\begin{proof}
\begin{align}
 e^{z/2}-1&=\tfrac z2\int_0^1e^{sz/2}\,ds,&
 \sum_{j=s}^{t}(u_{j+1}-u_j)&=u_{t+1}-u_s.
\end{align}
\end{proof}

\begin{lemma}[Logarithm and minimizing distance]
For $\alpha=-K\langle x,y\rangle_L\ge1$,
\begin{align}
 v&=\frac{\operatorname{arcosh}\alpha}{\sqrt{\alpha^2-1}}(y-\alpha x),\\
 \langle x,v\rangle_L&=0,&
 \|v\|^2&=K^{-1}(\operatorname{arcosh}\alpha)^2,&\Exp_x^Kv&=y.
\end{align}
For any piecewise $C^1$ curve $c:[0,1]\to H_K$ from $x$ to $y$,
\begin{align}
 \operatorname{Length}(c)&\ge K^{-1/2}\operatorname{arcosh}\alpha.
\end{align}
\end{lemma}
\begin{proof}
\begin{align}
 \sqrt{(K^{-1}+a^2)(K^{-1}+b^2)}&\ge K^{-1}+ab\quad(a,b\ge0),\\
 -\langle x,y\rangle_L&=x_0y_0-x_s^\top y_s\ge K^{-1},\\
 \langle y-\alpha x,y-\alpha x\rangle_L&=(\alpha^2-1)/K,\\
 \langle x,y-\alpha x\rangle_L&=-\alpha/K+\alpha/K=0,\\
 \Exp_x^Kv&=\alpha x+(y-\alpha x)=y.
\end{align}
For $\alpha(z)=-K\langle x,z\rangle_L$ and
$r_x(z)=K^{-1/2}\operatorname{arcosh}\alpha(z)$,
\begin{align}
 \operatorname{grad}r_x(z)
 &=\frac{-\sqrt K}{\sqrt{\alpha(z)^2-1}}\Pi_z^Kx,\\
 \|\Pi_z^Kx\|^2&=(\alpha(z)^2-1)/K,\\
 \|\operatorname{grad}r_x(z)\|&=1,\\
 r_x(y)-r_x(x)&=\int_0^1\langle\operatorname{grad}r_x(c(s)),\dot c(s)\rangle\,ds
 \le\int_0^1\|\dot c(s)\|\,ds.
\end{align}
The value at $z=x$ follows by continuity and the inequality by an arbitrarily short exclusion of that endpoint.
\end{proof}
